\chapter{Random walks and the Wiener process}
\label{ch:first-mg}

\begin{question}
Which processes are fair games? We meet the Wiener process, the most important process in this book, and find its three standard martingales.
\end{question}

\begin{whatfor}
\begin{itemize}
  \item Defines sample paths and the Wiener process (standard Brownian motion) by four rules.
  \item Proves that $W_t$, $W_t^2-t$ and $e^{uW_t-u^2t/2}$ are martingales.
  \item Finds when an additive or a geometric random walk is a martingale.
\end{itemize}
\skipline{fs:first-mg}
\end{whatfor}

\begin{prereq}
From Part~0: almost surely, independence, normal distribution, $\E_Pe^{u\xi}=e^{u\mu+u^2\sigma^2/2}$, sums of independent normals, CE1--CE4, Cauchy--Schwarz. From \cref{ch:martingales}: stochastic process, filtration, adapted, martingale, sub/supermartingale.
\end{prereq}

%=====================================================================
\section{Sample paths}

A stochastic process $(\xi(t))_{t\in\mathbb T}$ has two variables: time $t$ and outcome $\omega$. Fix $t$: you get a random variable $\xi(t,\cdot)$, as before. Fix $\omega$ instead:

\begin{definition}[Sample path]
For a fixed outcome $\omega$, the function $t\mapsto\xi(t,\omega)$ is a \dfn{sample path} (or path) of the process. The set of values the process can take is its \dfn{state space}; the set $\mathbb T$ of times is its \dfn{parameter space}.
\end{definition}

So a process is a random function: the experiment picks $\omega$, and with it a whole path.

\begin{pic}
\begin{tikzpicture}
\begin{axis}[width=10cm,height=5cm,xmin=0,xmax=1,ymin=-1.6,ymax=1.6,xlabel={$t$},ylabel={$W_t(\omega)$},legend style={font=\scriptsize,at={(0.02,0.98)},anchor=north west},legend columns=3]
  \addplot[defc,thick] table[x=t,y=w1] {figs/bm3.dat}; \addlegendentry{$\omega_1$}
  \addplot[thmc,thick] table[x=t,y=w2] {figs/bm3.dat}; \addlegendentry{$\omega_2$}
  \addplot[trapc,thick] table[x=t,y=w3] {figs/bm3.dat}; \addlegendentry{$\omega_3$}
  \draw[dashed,prec] (axis cs:0.6,-1.6)--(axis cs:0.6,1.6) node[pos=0.94,right,font=\scriptsize] {fix $t=0.6$};
\end{axis}
\end{tikzpicture}
\picnote{Three sample paths of one process (the process defined next), one for each of three outcomes. Reading across a path is one $\omega$; reading down the dashed line is one random variable $W_{0.6}$.}
\end{pic}

%=====================================================================
\section{The Wiener process}

\begin{definition}[Wiener process / standard Brownian motion]
\label{def:wiener}
A process $(W_t)_{t\ge0}$ is a \dfn{Wiener process}, or \dfn{standard Brownian motion} (SBM), if
\begin{enumerate}[label=(\arabic*)]
  \item $W_0=0$ a.s.;
  \item the sample paths $t\mapsto W_t$ are continuous a.s.;
  \item it has \dfn{independent increments}: for $0\le t_1<t_2<\dots<t_n$, the increments $W_{t_2}-W_{t_1},\dots,W_{t_n}-W_{t_{n-1}}$ are independent;
  \item for $0\le s<t$, the increment $W_t-W_s\sim N(0,\,t-s)$ (\dfn{stationary increments}: the law depends only on the length $t-s$).
\end{enumerate}
\end{definition}

That such a process exists is a theorem we take on trust. Chapter~8 shows where it comes from: it is the limit of a fair coin-tossing game with tiny steps.

Three facts follow at once.
\begin{itemize}
  \item $W_t=W_t-W_0\sim N(0,t)$: so $\E_P(W_t)=0$, $\E_P(W_t^2)=t$, $\E_Pe^{uW_t}=e^{u^2t/2}$.
  \item The future increment $W_t-W_s$ ($s<t$) is independent of the whole past $\F_s=\sigma(W_r:r\le s)$, the natural filtration. This is the property we use most.
  \item Integrable: by Cauchy--Schwarz, $\E_P|W_t|\le\sqrt{\E_PW_t^2}=\sqrt t<\infty$.
\end{itemize}

The one trick of this chapter is \emph{increment splitting}: write a future value as a known value plus an independent increment,
\[
  W_t=W_s+(W_t-W_s),\qquad\text{known at time $s$}\;+\;\text{independent of }\F_s,\ \sim N(0,t-s).
\]

%=====================================================================
\section{Three martingales of the Wiener process}

Throughout, $\F_t=\sigma(W_r:r\le t)$ and $0\le s<t$.

\begin{theorem}[The three canonical martingales]
\label{thm:canonical}
Each of the following is a martingale:
\begin{enumerate}[label=(\alph*)]
  \item $W_t$;
  \item $Y_t=W_t^2-t$;
  \item $Z_t=e^{uW_t-\frac12u^2t}$, for any number $u$; it has $\E_P(Z_t)=1$.
\end{enumerate}
\end{theorem}
\begin{proof}
All three are functions of $W_t$, so they are adapted. Split $W_t=W_s+(W_t-W_s)$ every time.
\begin{steps}
  \item (a) Integrable: $\E_P|W_t|\le\sqrt t$. By CE1, then CE2 and CE3,
  \[
    \E_P(W_t\given\F_s)=W_s+\E_P(W_t-W_s)=W_s+0=W_s .
  \]
  \item (b) Integrable: $\E_P|Y_t|\le\E_PW_t^2+t=2t$. Square the split, $W_t^2=W_s^2+2W_s(W_t-W_s)+(W_t-W_s)^2$, and condition term by term:
  \[
    \begin{aligned}
    \E_P(W_t^2\given\F_s)&=W_s^2+2W_s\,\E_P(W_t-W_s)+\E_P(W_t-W_s)^2\\
    &=W_s^2+0+(t-s).
    \end{aligned}
  \]
  Subtract $t$: $\E_P(W_t^2-t\given\F_s)=W_s^2-s=Y_s$.
  \item (c) $Z_t>0$ and $\E_PZ_t=e^{-u^2t/2}\E_Pe^{uW_t}=e^{-u^2t/2}e^{u^2t/2}=1$. Split the exponent:
  \[
    \begin{aligned}
    Z_t&=Z_s\cdot e^{u(W_t-W_s)-\frac12u^2(t-s)},\\
    \E_P(Z_t\given\F_s)&=Z_s\,e^{-\frac12u^2(t-s)}\,\E_Pe^{u(W_t-W_s)}=Z_s\,e^{-\frac12u^2(t-s)}\,e^{\frac12u^2(t-s)}=Z_s . \qedhere
    \end{aligned}
  \]
\end{steps}
\end{proof}

\begin{remark}
$W_t^2$ on its own is a \emph{sub}martingale: $\E_P(W_t^2\given\F_s)=W_s^2+(t-s)\ge W_s^2$. Subtracting $t$ removes exactly the upward drift. Likewise $e^{uW_t}$ is a submartingale, and the factor $e^{-u^2t/2}$ removes its drift.
\end{remark}

\begin{pic}
\begin{tikzpicture}
\begin{axis}[width=10cm,height=4.6cm,xmin=0,xmax=1,ymin=-1.3,ymax=1.6,xlabel={$t$},legend style={font=\scriptsize,at={(0.02,0.98)},anchor=north west},legend columns=2]
  \addplot[defc,thick] table[x=t,y expr=\thisrow{w1}^2] {figs/bm3.dat}; \addlegendentry{$W_t^2$ (drifts up)}
  \addplot[thmc,thick] table[x=t,y expr=\thisrow{w1}^2-\thisrow{t}] {figs/bm3.dat}; \addlegendentry{$W_t^2-t$ (fair)}
  \addplot[prec,dashed,domain=0:1] {x}; \addlegendentry{$\E_P W_t^2=t$}
  \addplot[prec,dotted,domain=0:1] {0}; \addlegendentry{$\E_P(W_t^2-t)=0$}
\end{axis}
\end{tikzpicture}
\picnote{One path of $W_t^2$ and of $W_t^2-t$. The mean of $W_t^2$ grows like $t$; subtracting $t$ leaves a process whose mean stays at $0$.}
\end{pic}

%=====================================================================
\section{Random walks}

\begin{definition}[Random walk; geometric random walk]
\label{def:rw}
Let $Y_1,Y_2,\ldots$ be i.i.d.\ (independent and identically distributed) random variables with $\E_P|Y_1|<\infty$, and $\F_n=\sigma(Y_1,\dots,Y_n)$.
\begin{itemize}
  \item The \dfn{random walk} is $X_0=0$, $X_n=Y_1+\dots+Y_n$.
  \item For a constant $X_0>0$, the \dfn{geometric random walk} is $X_n=X_0\,e^{Y_1+\dots+Y_n}$.
\end{itemize}
\end{definition}

\begin{result}[When is a walk a martingale?]
\label{res:rw-mg}
Split off the last step.
\begin{enumerate}[label=(\alph*)]
  \item $\E_P(X_{n+1}\given\F_n)=X_n+\E_P(Y_1)$ for the random walk, which is a martingale exactly when $\E_P(Y_1)=0$ (sub if $\ge0$, super if $\le0$).
  \item $\E_P(X_{n+1}\given\F_n)=X_n\,\E_P(e^{Y_1})$ for the geometric walk, which is a martingale exactly when $\E_P(e^{Y_1})=1$. If $Y_j\sim N(\mu,\sigma^2)$ this means $e^{\mu+\sigma^2/2}=1$, i.e.\ $\mu=-\frac12\sigma^2$.
\end{enumerate}
\end{result}
\begin{proof}
One step at a time.
\begin{steps}
  \item (a) $X_{n+1}=X_n+Y_{n+1}$; $X_n$ is known and $Y_{n+1}$ is independent of $\F_n$, so CE1--CE3 give $X_n+\E_P(Y_{n+1})=X_n+\E_P(Y_1)$. Integrable: $\E_P|X_n|\le n\,\E_P|Y_1|$.
  \item (b) $X_{n+1}=X_n\,e^{Y_{n+1}}$, so CE2--CE3 give $X_n\,\E_P(e^{Y_1})$. For normal $Y_1$ use $\E_Pe^{Y_1}=e^{\mu+\sigma^2/2}$. \qedhere
\end{steps}
\end{proof}

\begin{pic}
\begin{tikzpicture}
\begin{axis}[width=10cm,height=4.2cm,xmin=0,xmax=40,xlabel={$n$},ylabel={$X_n$},ymin=-8,ymax=8]
  \addplot[defc,thick,mark=*,mark size=0.8pt] table[x=k,y=s] {figs/walk.dat};
  \addplot[prec,dashed,domain=0:40] {0};
\end{axis}
\end{tikzpicture}
\picnote{A fair $\pm1$ random walk: each step adds $+1$ or $-1$ with probability $\frac12$. It is a martingale because each step has mean $0$.}
\end{pic}

\begin{example}[The drift correction]
\label{ex:drift-correction}
A stock modelled as $X_n=X_0e^{Y_1+\dots+Y_n}$ with daily log-returns $Y_j\sim N(\mu,\sigma^2)$, $\sigma=0.02$. It is a martingale only if $\mu=-\frac12(0.02)^2=-0.0002$. A zero median log-return ($\mu=0$) would make it a \emph{sub}martingale: $\E_Pe^{Y_1}=e^{0.0002}>1$. This $-\frac12\sigma^2$ is the same correction as in $e^{uW_t-\frac12u^2t}$; it returns in the stock-price models of Parts~2 and~4.
\end{example}

%=====================================================================
\formulasheet{fs:first-mg}

\begin{mastercard}[label=card:incr]{The three martingales of a process with independent increments}
$X$ has independent increments and starts at a fixed $X_0$. Split $X_t=X_s+(X_t-X_s)$ and use \cref{card:freeze}: take away what the fresh increment adds on average.
\[
  \boxed{\;X_t-\E_PX_t,\qquad \big(X_t-\E_PX_t\big)^2-\Var X_t,\qquad \frac{e^{uX_t}}{\E_P\,e^{uX_t}}\qquad\text{are martingales.}\;}
\]
\begin{plugrows}
  \plugrow{Standard Wiener process $W$}{\pstep{Box: $X_t-\E_PX_t$, $(X_t-\E_PX_t)^2-\Var X_t$, $e^{uX_t}/\E_Pe^{uX_t}$ are martingales.}\pstep{From $W$: $\E_PW_t=0$, $\Var W_t=t$.}\pstep{\cref{card:normal-mgf} with $\mu=0$, $\sigma^2=t$: $\E_Pe^{uW_t}=e^{u^2t/2}$.}\pstep{Substitute into the three.}}{$W_t$, $W_t^2-t$, $e^{uW_t-u^2t/2}$ are martingales}{$u$ any number}{$\E_Pe^{uW_t-u^2t/2}=e^{u^2t/2}\cdot e^{-u^2t/2}=1$ for all $t$}
  \plugrow{Conditional moments}{\pstep{Martingale (\cref{card:onestep}): $\E_P(M_t\given\F_s)=M_s$.}\pstep{$M=W$: $\E_P(W_t\given\F_s)=W_s$.}\pstep{$M=W^2-t$: $\E_P(W_t^2-t\given\F_s)=W_s^2-s$, so $\E_P(W_t^2\given\F_s)=W_s^2+(t-s)$.}\pstep{$M=e^{uW-u^2t/2}$: $\E_P(e^{uW_t}\given\F_s)e^{-u^2t/2}=e^{uW_s-u^2s/2}$, so $\E_P(e^{uW_t}\given\F_s)=e^{uW_s+u^2(t-s)/2}$.}}{$\E_P(W_t\given\F_s)=W_s$; $\E_P(W_t^2\given\F_s)=W_s^2+(t-s)$; $\E_P(e^{uW_t}\given\F_s)=e^{uW_s+u^2(t-s)/2}$}{$s<t$ (looking back, $s>t$, is Card 12.1)}{$\E_P(W_5^2\given\F_2)=W_2^2+(5-2)=W_2^2+3$; $u=2$: $\E_P(e^{2W_3}\given\F_1)=e^{2W_1+4\cdot2/2}=e^{2W_1+4}$}
  \plugrow{Sums of i.i.d.\ steps}{\pstep{Box with $X=S_n=Y_1+\dots+Y_n$.}\pstep{From i.i.d.\ steps: $\E_PS_n=n\mu$, $\Var S_n=n\sigma^2$ (\cref{card:var-lin}, row 4).}\pstep{\cref{card:indep}: $\E_Pe^{uS_n}=\prod\E_Pe^{uY_i}=(\E_Pe^{uY})^n$.}\pstep{Substitute into the three.}}{$S_n-n\mu$, $(S_n-n\mu)^2-n\sigma^2$, $e^{uS_n}/(\E_Pe^{uY})^n$ are martingales}{integer times}{$P(+1)=0.7$: $\mu=0.7-0.3=0.4$, so $S_n-0.4n$; fair $\pm1$: $\mu=0$, $\sigma^2=1$, so $S_n^2-n$; up $1$ w.p.\ $\frac67$, down $2$ w.p.\ $\frac17$, $e^u=\frac12$: $\E_P2^{-Y}=\frac12\cdot\frac67+4\cdot\frac17=\frac37+\frac47=1$, so $2^{-S_n}$}
  \plugrow{Exponential of a sum (geometric walk)}{\pstep{Row 3, third martingale, with $u=1$: $e^{S_n}/(\E_Pe^{Y})^n$.}\pstep{So $X_0e^{S_n}$ itself is a martingale iff $\E_Pe^Y=1$.}\pstep{For $Y\sim N(\mu,\sigma^2)$: \cref{card:normal-mgf}, $\E_Pe^Y=e^{\mu+\sigma^2/2}$, which is $1$ iff $\mu=-\frac12\sigma^2$.}}{$X_0e^{S_n}$ is a martingale iff $\E_Pe^{Y}=1$; for $Y\sim N(\mu,\sigma^2)$ iff $\mu=-\frac12\sigma^2$}{$\E_Pe^{Y}<\infty$}{$Y\sim N(\mu,0.04)$: $\mu=-\frac12(0.04)=-0.02$}
  \plugrow{Products of i.i.d.\ factors $\eta_i\ge0$}{\pstep{Row 4: $e^{S_n}/(\E_Pe^Y)^n$ is a martingale.}\pstep{Set $Y_i=\ln\eta_i$: $e^{S_n}=\eta_1\cdots\eta_n$ and $\E_Pe^{Y}=\E_P\eta$.}\pstep{Substitute.}}{$\prod_{i\le n}\eta_i\big/(\E_P\eta)^n$ is a martingale}{$\E_P\eta\in(0,\infty)$}{$U_i$ uniform, $\E_PU=\frac12$: $2^nU_1\cdots U_n$; $\eta=e^{-\epsilon}$, $\E_P\eta=\int_0^\infty e^{-x}e^{-x}dx=\frac12$: $2^ne^{-S_n}$; Bernoulli$(p)$, $\E_P\xi=p$: $p^{-n}\xi_1\cdots\xi_n$}
  \plugrow{Wiener process with drift and scale}{\pstep{Box with $X_t=x+\mu t+\sigma W_t$.}\pstep{From $X$: $\E_PX_t=x+\mu t$, $\Var X_t=\sigma^2t$.}\pstep{\cref{card:normal-mgf}: $\E_Pe^{uX_t}=e^{ux+(u\mu+\frac12u^2\sigma^2)t}$.}\pstep{Substitute; drop the constants $x$ and $e^{-ux}$ (a constant shift or factor keeps a martingale).}\pstep{Choose $u=-2\mu/\sigma^2$: then $u\mu+\frac12u^2\sigma^2=-\frac{2\mu^2}{\sigma^2}+\frac{2\mu^2}{\sigma^2}=0$.}}{$X_t-\mu t$, $(X_t-x-\mu t)^2-\sigma^2t$, $e^{uX_t-(u\mu+u^2\sigma^2/2)t}$; with $u=-2\mu/\sigma^2$: $e^{-2\mu X_t/\sigma^2}$}{$\sigma>0$}{$\mu=1.5$, $\sigma=1$: $u=-2(1.5)/1=-3$, so $e^{-3X_t}$}
  \plugrow{Poisson process, rate $\lambda$}{\pstep{Box with $X=N$.}\pstep{From Poisson$(\lambda t)$: $\E_PN_t=\Var N_t=\lambda t$.}\pstep{$\E_Pe^{uN_t}=\sum_ke^{uk}e^{-\lambda t}\frac{(\lambda t)^k}{k!}=e^{-\lambda t}e^{\lambda te^u}=e^{\lambda t(e^u-1)}$.}\pstep{Substitute into the three.}}{$N_t-\lambda t$, $(N_t-\lambda t)^2-\lambda t$, $e^{uN_t-\lambda t(e^u-1)}$ are martingales}{independent Poisson increments}{$N-\lambda t$ martingale: $\E_P(N_{t+s}-\lambda(t+s)\given\F_t)=N_t-\lambda t$, so $\E_P(N_{t+s}\given\F_t)=N_t+\lambda s$}
\end{plugrows}
\end{mastercard}
